\documentclass[11pt,a4paper]{article}
\usepackage[margin=2.5cm]{geometry}
\usepackage{hyperref}
\usepackage{listings}
\usepackage{xcolor}
\usepackage{booktabs}
\usepackage{longtable}
\usepackage{array}
\usepackage{parskip}
\usepackage{titlesec}
\usepackage{fancyhdr}
\usepackage{microtype}

% ── Code listing style ──────────────────────────────────────────────────────
\lstdefinelanguage{TextIR}{
  morekeywords={module,version,class,interface,struct,enum,field,method,
                constructor,local,public,private,protected,internal,static,
                virtual,override,abstract,sealed,implements,else,catch,finally},
  morestring=[b]",
  morecomment=[l]{//},
  sensitive=true
}
\lstdefinestyle{code}{
  basicstyle=\ttfamily\small,
  keywordstyle=\color{blue}\bfseries,
  commentstyle=\color{gray}\itshape,
  stringstyle=\color{teal},
  frame=single,
  rulecolor=\color{black!20},
  backgroundcolor=\color{black!3},
  breaklines=true,
  numbers=left,
  numberstyle=\tiny\color{gray},
  numbersep=6pt,
  xleftmargin=1.5em
}
\lstset{style=code}

% ── Page style ───────────────────────────────────────────────────────────────
\pagestyle{fancy}
\fancyhf{}
\rhead{Lattice IRRuntime}
\lhead{Feature Reference}
\rfoot{\thepage}

% ── Title ───────────────────────────────────────────────────────────────────
\title{\textbf{Lattice IRRuntime}\\\large Feature Set Reference}
\author{Lattice Project}
\date{March 2026}

\begin{document}

\maketitle
\tableofcontents
\newpage

% ============================================================================
\section{Overview}
% ============================================================================

\textbf{Lattice} is a stack-based virtual machine written in C\# (.NET 10).
It executes programs described in an \emph{Intermediate Representation} (IR).
The runtime loads an IR module, resolves types and methods, and executes
instructions on a per-frame evaluation stack.

Key design goals:
\begin{itemize}
  \item Simple, human-readable text IR format (\texttt{TextIR}).
  \item Seamless interop with native C\# types and methods.
  \item Embeddable --- host applications can load a module, call IR methods,
        and register C\# types or callbacks without writing a parser.
\end{itemize}

% ============================================================================
\section{Input Formats}
% ============================================================================

The runtime accepts three equivalent input representations, selected via
the \texttt{InputFormat} enum or auto-detected.

\begin{center}
\begin{tabular}{lll}
\toprule
\textbf{Format} & \textbf{Enum value} & \textbf{Description} \\
\midrule
\texttt{Auto}   & \texttt{InputFormat.Auto}   & JSON if starts with \texttt{\{} or \texttt{[}, else TextIR \\
\texttt{TextIr} & \texttt{InputFormat.TextIr} & Human-readable text IR (see Section~\ref{sec:textir}) \\
\texttt{Json}   & \texttt{InputFormat.Json}   & Serialised \texttt{ModuleDto} JSON \\
\texttt{Ast}    & \texttt{InputFormat.Ast}    & \texttt{ModuleNode} AST lowered by \texttt{AstLowering} \\
\bottomrule
\end{tabular}
\end{center}

% ============================================================================
\section{TextIR Syntax}
\label{sec:textir}
% ============================================================================

A TextIR file begins with a module header, followed by one or more type
definitions.

\begin{lstlisting}[language=TextIR,caption={Minimal TextIR module}]
module MyApp version 1.0.0

class Program {
    static method Main() -> void {
        ldstr "Hello, World!"
        call System.Console.WriteLine(System.string) -> void
        ret
    }
}
\end{lstlisting}

\subsection{Module header}
\begin{lstlisting}[language=TextIR]
module <name> version <major>.<minor>[.<patch>]
\end{lstlisting}

\subsection{Type declaration keywords}
\begin{center}
\begin{tabular}{ll}
\toprule
\textbf{Keyword} & \textbf{Meaning} \\
\midrule
\texttt{class}     & Reference type \\
\texttt{interface} & Interface (pure virtual contract) \\
\texttt{struct}    & Value type \\
\texttt{enum}      & Enumeration type \\
\bottomrule
\end{tabular}
\end{center}

Type modifiers (\texttt{public}, \texttt{private}, \texttt{abstract},
\texttt{sealed}) may precede the keyword.  Inheritance is expressed with
\texttt{:}:

\begin{lstlisting}[language=TextIR]
class Animal {
    field name: string
    method Speak() -> void { ... }
}

class Dog : Animal {
    override method Speak() -> void implements Animal.Speak { ... }
}
\end{lstlisting}

\subsection{Member declarations}

\begin{lstlisting}[language=TextIR]
// Field
[static] field <name>: <type>

// Method
[static] [virtual] [override] [abstract] method <name>(<params>) -> <rettype> [implements <Interface>.<MethodName>] { ... }

// Constructor
constructor(<params>) { ... }

// Local variable
local <name>: <type>
\end{lstlisting}

% ============================================================================
\section{Type System}
% ============================================================================

\subsection{Primitive types}

\begin{longtable}{lll}
\toprule
\textbf{IR name(s)} & \textbf{Alias(es)} & \textbf{CLR type} \\
\midrule
\texttt{void}   & \texttt{system.void}    & \texttt{void} \\
\texttt{bool}   & \texttt{system.boolean} & \texttt{bool} \\
\texttt{int8}   & \texttt{system.sbyte}   & \texttt{sbyte} \\
\texttt{uint8}  & \texttt{system.byte}    & \texttt{byte} \\
\texttt{int16}  & \texttt{system.int16}   & \texttt{short} \\
\texttt{uint16} & \texttt{system.uint16}  & \texttt{ushort} \\
\texttt{int32}  & \texttt{system.int32}   & \texttt{int} \\
\texttt{uint32} & \texttt{system.uint32}  & \texttt{uint} \\
\texttt{int64}  & \texttt{system.int64}   & \texttt{long} \\
\texttt{uint64} & \texttt{system.uint64}  & \texttt{ulong} \\
\texttt{float32}, \texttt{single} & \texttt{system.single} & \texttt{float} \\
\texttt{float64}, \texttt{double} & \texttt{system.double} & \texttt{double} \\
\texttt{char}   & \texttt{system.char}    & \texttt{char} \\
\texttt{string} & \texttt{system.string}  & \texttt{string} \\
\texttt{object} & \texttt{system.object}  & \texttt{object} \\
\texttt{decimal}  & \texttt{system.decimal}  & \texttt{decimal} \\
\texttt{datetime} & \texttt{system.datetime} & \texttt{DateTime} \\
\texttt{timespan} & \texttt{system.timespan} & \texttt{TimeSpan} \\
\texttt{guid}     & \texttt{system.guid}     & \texttt{Guid} \\
\bottomrule
\end{longtable}

Type names are case-insensitive in IR operands.  The runtime normalises them
to lower-case before lookup, so \texttt{System.String}, \texttt{system.string},
and \texttt{string} are all equivalent.

\subsection{User-defined types}

User-defined types are declared in the module and managed by the VM as
\texttt{ManagedObject} instances at runtime.  Each instance carries:
\begin{itemize}
  \item A type name string (\texttt{TypeName}).
  \item A dictionary of named field values.
  \item Optional CLR-backing object (see Section~\ref{sec:clrtype}).
  \item Attached attribute instances.
  \item Arbitrary metadata slots.
\end{itemize}

% ============================================================================
\section{Instruction Set}
% ============================================================================

\subsection{Stack manipulation}

\begin{center}
\begin{tabular}{lp{10cm}}
\toprule
\textbf{Opcode} & \textbf{Description} \\
\midrule
\texttt{nop}    & No operation. \\
\texttt{dup}    & Duplicate the top value on the evaluation stack. \\
\texttt{pop}    & Discard the top value. \\
\texttt{ldnull} & Push \texttt{null} onto the stack. \\
\bottomrule
\end{tabular}
\end{center}

\subsection{Load constants}

\begin{center}
\begin{tabular}{lp{9cm}}
\toprule
\textbf{Opcode} & \textbf{Description} \\
\midrule
\texttt{ldc}      & Push a typed constant.  Operand: \texttt{\{value, type\}}.  Type defaults to \texttt{int64}. \\
\texttt{ldc.i4}   & Push a 32-bit integer constant. \\
\texttt{ldc.i8}   & Push a 64-bit integer constant. \\
\texttt{ldc.r4}   & Push a 32-bit float constant. \\
\texttt{ldc.r8}   & Push a 64-bit float constant. \\
\texttt{ldstr}    & Push a string constant.  Operand: \texttt{\{value\}}. \\
\bottomrule
\end{tabular}
\end{center}

\subsection{Variables and arguments}

\begin{center}
\begin{tabular}{lp{9cm}}
\toprule
\textbf{Opcode} & \textbf{Description} \\
\midrule
\texttt{ldarg}  & Push a method argument by index or name (including \texttt{this}). \\
\texttt{starg}  & Pop and store into a named argument slot. \\
\texttt{ldloc}  & Push a local variable by name. \\
\texttt{stloc}  & Pop and store into a local variable by name. \\
\bottomrule
\end{tabular}
\end{center}

\subsection{Fields}

\begin{center}
\begin{tabular}{lp{9cm}}
\toprule
\textbf{Opcode} & \textbf{Description} \\
\midrule
\texttt{ldfld}  & Pop an object from the stack (or use \texttt{this}) and push the named instance field. \\
\texttt{stfld}  & Pop a value, then pop an object; store the value into the named instance field. \\
\texttt{ldsfld} & Push a static field.  Operand: \texttt{\{field: \{declaringType, name\}\}}. \\
\texttt{stsfld} & Pop a value; store it into a static field. \\
\bottomrule
\end{tabular}
\end{center}

\noindent When the object wraps a CLR-backed instance
(see Section~\ref{sec:clrtype}), \texttt{ldfld}/\texttt{stfld} use
reflection to read/write the corresponding CLR property or field.

\subsection{Arithmetic}

\begin{center}
\begin{tabular}{lp{9cm}}
\toprule
\textbf{Opcode} & \textbf{Description} \\
\midrule
\texttt{add}  & Pop $b$, pop $a$, push $a + b$. \\
\texttt{sub}  & Push $a - b$. \\
\texttt{mul}  & Push $a \times b$. \\
\texttt{div}  & Push $a \div b$. \\
\texttt{rem}  & Push $a \bmod b$. \\
\texttt{neg}  & Pop $v$, push $-v$. \\
\bottomrule
\end{tabular}
\end{center}

If either operand is a floating-point value, both are promoted to
\texttt{double}.  Integer operations return \texttt{long}.

\subsection{Logical}

\begin{center}
\begin{tabular}{lp{9cm}}
\toprule
\textbf{Opcode} & \textbf{Description} \\
\midrule
\texttt{not} & Pop $v$, push \texttt{!ToBool($v$)}. \\
\bottomrule
\end{tabular}
\end{center}

\subsection{Comparison}

Each opcode pops $b$ then $a$ and pushes a \texttt{bool}.

\begin{center}
\begin{tabular}{ll}
\toprule
\textbf{Opcode} & \textbf{Condition} \\
\midrule
\texttt{ceq} & $a = b$ \\
\texttt{cne} & $a \neq b$ \\
\texttt{cgt} & $a > b$ \\
\texttt{cge} & $a \geq b$ \\
\texttt{clt} & $a < b$ \\
\texttt{cle} & $a \leq b$ \\
\bottomrule
\end{tabular}
\end{center}

If either operand is a \texttt{string} the comparison is performed
lexicographically (ordinal).  Boolean operands are supported for
\texttt{ceq}/\texttt{cne}.  All other operands are promoted to
\texttt{double}.

\subsection{Objects and arrays}

\begin{center}
\begin{tabular}{lp{9cm}}
\toprule
\textbf{Opcode} & \textbf{Description} \\
\midrule
\texttt{newobj}   & Create a new \texttt{ManagedObject} of the given type.  If the type is registered as a CLR type, a real CLR instance is created and attached. \\
\texttt{newarr}   & Push a new \texttt{List<object?>} (dynamically-typed array). \\
\texttt{ldelem}   & Pop index, pop array, push \texttt{array[index]}. \\
\texttt{stelem}   & Pop value, pop index, pop array; set \texttt{array[index] = value}. \\
\bottomrule
\end{tabular}
\end{center}

\subsection{Type conversion and testing}

\begin{center}
\begin{tabular}{lp{9cm}}
\toprule
\textbf{Opcode} & \textbf{Description} \\
\midrule
\texttt{conv}       & Pop a value and push it converted to the named target type. \\
\texttt{castclass}  & Assert that the top-of-stack \texttt{ManagedObject} has the expected type name; throws \texttt{InvalidCastException} otherwise. \\
\texttt{isinst}     & Pop a value; push \texttt{true} if it is a \texttt{ManagedObject} with the given type name, otherwise \texttt{false}. \\
\bottomrule
\end{tabular}
\end{center}

\subsection{Method calls}

\begin{center}
\begin{tabular}{lp{9cm}}
\toprule
\textbf{Opcode} & \textbf{Description} \\
\midrule
\texttt{call}     & Static or instance call.  Operand: \texttt{\{method: \{declaringType, name, returnType, parameterTypes[]\}\}}. \\
\texttt{callvirt} & Virtual call.  Identical to \texttt{call} but additionally pops an object instance from the stack. \\
\texttt{ret}      & Return from the current frame.  If a value is on the stack it is returned to the caller. \\
\bottomrule
\end{tabular}
\end{center}

\paragraph{Call resolution order}
\begin{enumerate}
  \item Exact signature lookup in the native method dictionary.
  \item Simplified-signature lookup (parameter types reduced to unqualified
        lower-case short names, e.g.\ \texttt{System.String} $\to$ \texttt{string}).
  \item Reflection-based CLR invocation (only when
        \texttt{enableReflectionNativeMethods} is \texttt{true}).
  \item IR method lookup in the loaded module.
\end{enumerate}

\subsection{Control flow}

\begin{center}
\begin{tabular}{lp{9cm}}
\toprule
\textbf{Opcode}    & \textbf{Description} \\
\midrule
\texttt{if}        & Pop (or evaluate) a condition; execute \texttt{thenBlock} or optional \texttt{elseBlock}. \\
\texttt{while}     & Repeatedly evaluate condition and execute \texttt{body}.  Supports \texttt{break}/\texttt{continue}. \\
\texttt{break}     & Exit the nearest enclosing \texttt{while} loop. \\
\texttt{continue}  & Skip to the next iteration of the nearest \texttt{while} loop. \\
\texttt{try}       & Structured exception handling with \texttt{tryBlock}, zero or more \texttt{catchBlocks}, and optional \texttt{finallyBlock}. \\
\texttt{throw}     & Pop a value (VM exception object or CLR exception) and set it as a pending exception on the current frame. \\
\bottomrule
\end{tabular}
\end{center}

\paragraph{Condition kinds} A condition embedded in \texttt{if} or
\texttt{while} is one of:

\begin{center}
\begin{tabular}{ll}
\toprule
\textbf{Kind}      & \textbf{Evaluation} \\
\midrule
\texttt{stack}      & Pop the top-of-stack value and coerce to \texttt{bool}. \\
\texttt{binary}     & Pop $r$, pop $l$; compare using the specified \texttt{operation} opcode. \\
\texttt{expression} & Execute a single embedded \texttt{InstructionDto} and pop the result. \\
\texttt{block}      & Execute a sequence of instructions and pop the final result. \\
\bottomrule
\end{tabular}
\end{center}

% ============================================================================
\section{Runtime API}
% ============================================================================

\subsection{Construction}

\begin{lstlisting}[language={[Sharp]C}]
// Load from a TextIR/JSON/AST string at construction time
var rt = new IRRuntime(irSource);

// Construct empty, load later
var rt = new IRRuntime();
rt.LoadModule(irSource, InputFormat.TextIr);

// Enable reflection-based CLR interop
var rt = new IRRuntime(irSource, enableReflectionNativeMethods: true);
\end{lstlisting}

\subsection{Running a program}

\begin{lstlisting}[language={[Sharp]C}]
// Requires a static Program.Main() method in the module
rt.Run();
\end{lstlisting}

\subsection{Calling individual methods}

\begin{lstlisting}[language={[Sharp]C}]
// Untyped call
object? result = rt.CallMethod("MyType.MyMethod", arg1, arg2);

// Generic typed call with auto-conversion
int n = rt.CallMethod<int>("Calculator.Add", 3, 4);
\end{lstlisting}

\subsection{Registering native methods}

Bind a C\# delegate to an IR call signature:

\begin{lstlisting}[language={[Sharp]C}]
rt.RegisterNativeMethod(
    "MathHelper.Sqrt(float64)",
    (instance, args) => Math.Sqrt((double)args[0]!));
\end{lstlisting}

The signature format is:
\[
\texttt{DeclaringType.MethodName(param1Type,param2Type,...)}
\]
Both fully-qualified (\texttt{System.String}) and short (\texttt{string})
parameter type names are tried automatically.

\subsection{Observability hooks}

\begin{lstlisting}[language={[Sharp]C}]
// Called before every instruction (debugger / profiler hook)
rt.OnStep = (frame, instr) => Console.WriteLine(instr.opCode);

// Called when an unhandled exception escapes
rt.OnException = ex => Console.Error.WriteLine(ex);
\end{lstlisting}

\subsection{Accessing the call stack}

\begin{lstlisting}[language={[Sharp]C}]
IReadOnlyList<CallFrame> stack = rt.CallStack;
\end{lstlisting}

% ============================================================================
\section{CLR Type Registration}
\label{sec:clrtype}
% ============================================================================

Host code can register C\# classes so the VM can instantiate, field-access,
and call methods on them as if they were IR-defined types.

\begin{lstlisting}[language={[Sharp]C}]
// Register by type parameter (uses FullName as IR name by default)
IRRuntime.RegisterClrType<MyService>();

// Register with a custom IR name
IRRuntime.RegisterClrType<MyService>("Services.MyService");

// Register via a Type object
IRRuntime.RegisterClrType("Services.MyService", typeof(MyService));
\end{lstlisting}

Once registered:
\begin{itemize}
  \item \textbf{\texttt{newobj}} creates a real CLR instance via
        \texttt{Activator.CreateInstance} and attaches it to the
        \texttt{ManagedObject} wrapper.
  \item \textbf{\texttt{ldfld}/\texttt{stfld}} use reflection to
        read/write properties and fields on the CLR instance, falling
        back to the VM field dictionary if the member is not found.
  \item \textbf{\texttt{ldsfld}/\texttt{stsfld}} reflectively access
        static CLR fields/properties.
  \item \textbf{Reflected calls} (\texttt{EnableReflectionNativeMethods})
        receive the real CLR instance, not the \texttt{ManagedObject} wrapper.
  \item The type name is also consulted by \texttt{ResolveClrType}, so it
        can be used in cast expressions and catch clauses.
\end{itemize}

\paragraph{Example: Registering a custom service}
\begin{lstlisting}[language={[Sharp]C}]
public class Counter {
    public int Value { get; set; }
    public void Increment() => Value++;
    public int GetValue() => Value;
}

IRRuntime.RegisterClrType<Counter>("Counter");

var rt = new IRRuntime("""
    module CounterDemo version 1.0

    class Program {
        static method Main() -> void {
            newobj Counter
            callvirt Counter.Increment()
            callvirt Counter.GetValue() -> int32
            call System.Console.WriteLine(int32) -> void
            ret
        }
    }
""", enableReflectionNativeMethods: true);

rt.Run(); // prints: 1
\end{lstlisting}

% ============================================================================
\section{Built-in Native Methods}
% ============================================================================

The following native methods are pre-registered on every \texttt{IRRuntime}
instance.

\begin{longtable}{lp{8cm}}
\toprule
\textbf{Signature} & \textbf{Behaviour} \\
\midrule
\texttt{System.Console.WriteLine(string)}      & \texttt{Console.WriteLine(args[0])} \\
\texttt{System.Console.WriteLine(int32)}       & \texttt{Console.WriteLine(args[0])} \\
\texttt{System.Console.WriteLine(object)}      & \texttt{Console.WriteLine(args[0])} \\
\texttt{System.Console.WriteLine(string,string)} & Concatenates all args and writes a line \\
\texttt{System.Console.WriteLine(string,object)} & \texttt{string.Format} then writes a line \\
\texttt{System.Console.Write(string)}          & \texttt{Console.Write(args[0])} \\
\texttt{System.Console.Write(string,string)}   & Concatenates all args and writes \\
\texttt{System.Console.ReadLine()}             & Returns \texttt{Console.ReadLine()} \\
\texttt{System.Console.Clear()}                & \texttt{Console.Clear()} \\
\texttt{System.Console.Beep()}                 & \texttt{Console.Beep()} \\
\texttt{System.String.Concat(string,string)}   & \texttt{string.Concat(args[0], args[1])} \\
\texttt{System.String.Format(string,object)}   & \texttt{string.Format(fmt, args[1])} \\
\texttt{System.String.Format(string,object,object)} & \texttt{string.Format(fmt, a, b)} \\
\texttt{System.String.IsNullOrEmpty(string)}   & \texttt{string.IsNullOrEmpty(args[0])} \\
\texttt{System.Object.ToString()}              & \texttt{instance?.ToString()} \\
\texttt{System.Object.GetType()}               & \texttt{instance?.GetType()} \\
\bottomrule
\end{longtable}

% ============================================================================
\section{Value Coercion Rules}
% ============================================================================

\begin{center}
\begin{tabular}{lp{9cm}}
\toprule
\textbf{Target} & \textbf{Conversion} \\
\midrule
\texttt{bool}   & \texttt{null}/\texttt{0}/\texttt{""} $\to$ \texttt{false}; everything else $\to$ \texttt{true}. \\
\texttt{long}   & Numeric types truncated/cast; strings parsed with \texttt{InvariantCulture}. \\
\texttt{double} & Same as \texttt{long} but without truncation. \\
\texttt{string} & \texttt{object.ToString()}, or \texttt{null} if value is \texttt{null}. \\
CLR type        & \texttt{Convert.ChangeType} fallback after the typed fast-paths. \\
\bottomrule
\end{tabular}
\end{center}

% ============================================================================
\section{Exception Model}
% ============================================================================

\begin{itemize}
  \item IR code raises exceptions with \texttt{throw}, which pops a value
        (any object, including a \texttt{ManagedObject}) and stores it as
        \texttt{frame.PendingException}.
  \item \texttt{try}/\texttt{catch}/\texttt{finally} blocks are first-class
        instructions handled inline by the interpreter.
  \item Catch clauses match by:
        \begin{enumerate}
          \item Catch-all (no type specified).
          \item IR \texttt{ManagedObject} type name match.
          \item CLR \texttt{IsAssignableFrom} check for CLR exception types.
        \end{enumerate}
  \item If no catch clause matches, the pending exception propagates up the
        call stack.  At the top level it is delivered to \texttt{OnException}.
\end{itemize}

% ============================================================================
\section{Attributes}
% ============================================================================

Custom attributes can be applied to IR types, fields, and methods using
the standard \texttt{@AttributeName(args)} syntax in the TextIR.  At
runtime, attribute instances are created as \texttt{ManagedObject}s and
attached to the corresponding \texttt{ManagedObject} via
\texttt{AttachAttributesFromIR}.  Host code can retrieve them with:

\begin{lstlisting}[language={[Sharp]C}]
ManagedObject instance = ...;
var attr = instance.GetAttribute<MyManagedAttribute>();
\end{lstlisting}

% ============================================================================
\section{Limitations and Known Gaps}
% ============================================================================

\begin{itemize}
  \item Generics and generic type parameters are not yet supported.
  \item Arrays are implemented as \texttt{List<object?>}; typed arrays
        are not distinguished at runtime.
  \item \texttt{castclass} and \texttt{isinst} only match by exact
        \texttt{TypeName} string; inheritance-aware checks are not
        performed for IR types.
  \item \texttt{newobj} uses \texttt{Activator.CreateInstance()};
        constructor arguments for CLR-backed types must be passed via
        native method wrappers or IR constructors.
  \item The \texttt{while (stack)} condition idiom requires the caller to
        push a fresh truthy value before each iteration; an exhausted
        stack is treated as \texttt{false} (loop exits gracefully).
\end{itemize}

\end{document}
